\section{The literal fixed-shift residual-to-native interface}
\label{sec:tpc58-fixed-shift-interface}

We now identify the high residual vector in
\cref{sec:tpc58-two-gram-dictionary} with the literal terminal output of
TPC--53--57.  This is a coordinatewise interface theorem at the one
specified \(h_0\).  It is not a lower Gram estimate.

\subsection{One side, one source system, and one physical intercept}

Fix one side of the equality column.  On the left, a terminal atom is
\begin{equation}
 \mathfrak u=(m,r,j,\gamma),
 \qquad m=d_m\ell_m,
 \qquad p_{\mathfrak u}r-mj=h_0,
 \label{eq:tpc58-left-terminal-atom}
\end{equation}
and its native and residual labels are
\begin{equation}
 i(\mathfrak u):=(L,\gamma,j),
 \qquad s(\mathfrak u):=d_mr.
 \label{eq:tpc58-left-native-residual-labels}
\end{equation}
The right side has
\(\mathfrak u=(m,r,j,\alpha)\), with the source systems interchanged, and
\(i(\mathfrak u)=(R,\alpha,j)\).  Every assertion below is first proved on
one side.  The two-sided assertion is the orthogonal direct sum of the two
one-sided maps; it is not a new coherent sum.

The row, mask, and scale data are the literal ones from TPC--53:
\begin{align}
 m&=d_m\ell_m,
 &a_m&=\mu(d_m)(\log\ell_m)\omega_D(d_m)
       \psi_L(\ell_m/L)\zeta_m,
 \label{eq:tpc58-literal-source-row}\\
 K_{m\gamma}
 &=\one_{\ell_m\ne\ell_\gamma}
   \one_{|m-m_\gamma|>\Delta}
   \one_{(d_m,d_\gamma)\le G},
 &\Delta&=QX^{-1/400},\qquad G=X^{1/400}.
 \label{eq:tpc58-literal-mask-and-scales}
\end{align}
The fixed opposite-row factor, orbit scalar, and four-factor profile are
also retained.  We abbreviate their exact product by
\begin{equation}
 \mathfrak P_{m,\gamma}(j)
 :=\eta_\gamma\mathscr R(q_\gamma)\xi_a
   W_1(mj/X)\Psi_1(d_mj/K)
   W_2(m_\gamma j/X)\Psi_2(d_\gamma j/K).
 \label{eq:tpc58-literal-profile-product}
\end{equation}
Thus the pre-terminal coordinate is
\begin{equation}
 [G_X^L]_{\mathfrak u}
 =a_mK_{m\gamma}\mathfrak P_{m,\gamma}(j).
 \label{eq:tpc58-literal-preterminal-coordinate}
\end{equation}

Use the normalized shaped terminal weight
\begin{equation}
 \Lambda_X(n):=(\log n)\Phi_0(n/X).
 \label{eq:tpc58-normalized-terminal-weight}
\end{equation}
This is the same inherited cutoff as the TPC--45 notation
\(\Lambda(n)=(\log n)\Phi_0(n)\) after its scale has been absorbed into
the profile.  We use \eqref{eq:tpc58-normalized-terminal-weight} throughout
and do not identify it with the von Mangoldt function.

Let \(\cU_{h_0}^{\rm term}\) be the exact terminal-admissible set: the
candidate \(p_{\mathfrak u}=(mj+h_0)/r\) is prime in the inherited support,
\(\Lambda_X(p_{\mathfrak u}r)\ne0\), and
\begin{equation}
 d_m,r,s=d_mr\ \text{are squarefree},
 \qquad(d_m,r)=1.
 \label{eq:tpc58-squarefree-terminal-support}
\end{equation}
For comparison with the TPC--45 very-high Gram, make the deterministic
restriction
\begin{equation}
 \cU_{h_0,>y}^{\rm term}
 :=\{\mathfrak u\in\cU_{h_0}^{\rm term}:p_{\mathfrak u}>y\},
 \label{eq:tpc58-very-high-terminal-set}
\end{equation}
where \(y\) obeys
\eqref{eq:tpc58-cutoff-geometric-condition}--
\eqref{eq:tpc58-cutoff-label-separation}.  The conclusions in this section
concern this very-high face.  They do not silently cover the remainder of
the large-prime band.  At the endpoint cutoff, the cofactor ceiling
\eqref{eq:tpc58-compatible-cofactor-ceiling} also makes this face disjoint
from the old \(R\asymp J\) block.  Nontrivial activation on a compatible
smaller cofactor band is an additional arithmetic hypothesis, not a
consequence of the coordinate dictionary.

\subsection{Baseline and distinguished supports}

The exact factorization of TPC--56--57 has the form
\begin{equation}
 X_{\mathfrak u}
 :=(D_{h_0}G_X^L)_{\mathfrak u}
 =c_{\mathfrak u}\zeta_{m(\mathfrak u)},
 \label{eq:tpc58-terminal-row-factorization}
\end{equation}
where \(c_{\mathfrak u}\) contains the common output factor, literal mask,
all displayed profiles, source shapes, and \(\Lambda_X(mj+h_0)\), but not
the bounded row coefficient \(\zeta_m\).  Define the baseline and physical
very-high carriers separately:
\begin{align}
 \cU_{h_0,>y}^{\circ}
 &:=\{\mathfrak u\in\cU_{h_0,>y}^{\rm term}:
                 c_{\mathfrak u}\ne0\},
 \label{eq:tpc58-baseline-terminal-carrier}\\
 \cU_{h_0,>y}^{\rm phys}(\zeta)
 &:=\{\mathfrak u\in\cU_{h_0,>y}^{\circ}:
                 \zeta_{m(\mathfrak u)}\ne0\}.
 \label{eq:tpc58-physical-terminal-carrier}
\end{align}
The Hilbert spaces and the erasure map are built on the fixed baseline
carrier \eqref{eq:tpc58-baseline-terminal-carrier}.  A distinguished
coefficient vector is represented by zeros on inactive baseline
coordinates.  This keeps coefficient-dependent cancellation separate from
structural deletion.

TPC--55 proves, under \(D=o(J)\), row-factorization uniqueness, and large
\(X\), that
\begin{equation}
 \mathfrak u\longmapsto
 \bigl(i(\mathfrak u),s(\mathfrak u),p_{\mathfrak u}\bigr)
 \label{eq:tpc58-prime-resolved-injection}
\end{equation}
is injective on one declared source system.  Thus a coordinate
\((i,s,p)\) denotes at most one literal atom.  This is the canonical
nonduplicated prime-resolved representation on the fixed baseline carrier
used below; it is not a claim of minimal parameter dimension or physical
identifiability \citep{WangTPC55}.

\subsection{Sign calibration against the stripped TPC--45 atom}

Let \(b_{i,sp}\) be the stripped TPC--45 coefficient on the coordinate
corresponding through
\eqref{eq:tpc58-prime-resolved-injection}.  The exact factor allocation of
TPC--53 puts \(\mu(d_m)\) in the source coefficient and the terminal
operator contributes \(\mu(r)\).  Hence
\begin{equation}
 \mu(d_m)\mu(r)=\mu(s)
 \label{eq:tpc58-common-mobius-product}
\end{equation}
on \eqref{eq:tpc58-squarefree-terminal-support}.  With the leading sign in
the TPC--45 stripped atom fixed as in that paper, coordinatewise comparison
of the source coefficient, mask, opposite-row factor, orbit scalar,
profiles, and shaped terminal weight gives
\begin{equation}
 \boxed{
 X_{\mathfrak u}=-\mu(s(\mathfrak u))
 b_{i(\mathfrak u),\,s(\mathfrak u)p_{\mathfrak u}}.}
 \label{eq:tpc58-literal-sign-calibration}
\end{equation}
No modulus has been taken in
\eqref{eq:tpc58-literal-sign-calibration}.  It is the coefficientwise
global-sign calibration needed for the later base--high affine cross term
\citep{WangTPC45,WangTPC53}.

For clarity, define the stripped grouped high vector directly from the
TPC--45 atoms by
\begin{equation}
 z_H(i,s)
 :=-\sum_{\substack{p>y\\sp\in\cU_i^\circ}}b_{i,sp}.
 \label{eq:tpc58-literal-stripped-high-vector}
\end{equation}
This is exactly the convention \(z_H=-\sum_p x_p\) of
\eqref{eq:tpc58-high-residual-sign-convention}.  Put
\begin{equation}
 A_H:=\Uop_\mu z_H,
 \qquad
 A_H(i,s)=\mu(s)z_H(i,s).
 \label{eq:tpc58-literal-physical-high-vector}
\end{equation}
Then \eqref{eq:tpc58-literal-sign-calibration} and the prime-resolved
injection show
\begin{equation}
 \boxed{
 A_H(i,s)
 =\sum_{\substack{\mathfrak u\in\cU_{h_0,>y}^{\rm term}\\
                   i(\mathfrak u)=i,\ s(\mathfrak u)=s}}
   X_{\mathfrak u}.}
 \label{eq:tpc58-terminal-grouped-equals-physical-high}
\end{equation}
Thus \(A_H\), not the stripped \(z_H\), is the TPC--57 terminal grouped
amplitude.  In particular, the TPC--57 amplitude already contains the
common M\"obius factor \(\mu(s)\); applying \(\Uop_\mu\) to it a second time
would remove rather than restore the physical residual sign.

\begin{proposition}[Literal fixed-shift erasure interface]
\label{prop:tpc58-literal-fixed-shift-erasure}
On one declared side and source system, under the cutoff conditions
\eqref{eq:tpc58-cutoff-geometric-condition}--
\eqref{eq:tpc58-cutoff-label-separation} and the TPC--55 injection,
the following diagram commutes:
\begin{equation}
 \begin{array}{ccc}
 \{\text{prime-resolved atoms }(i,s,p)\}
 &\xrightarrow{\quad p\text{-grouping}\quad}&
 A_H\in\ell^2(i,s)\\[3pt]
 &&\ \downarrow\Eop\\[-1pt]
 &&-\displaystyle\sum_{p\in\cP_y}v_p\in\ell^2(i).
 \end{array}
 \label{eq:tpc58-fixed-shift-commuting-diagram}
\end{equation}
Equivalently,
\begin{equation}
 \boxed{
 \Eop A_H=-\sum_{p\in\cP_y}v_p.}
 \label{eq:tpc58-fixed-shift-native-high-identity}
\end{equation}
Every arrow is pointwise at the one physical \(h_0\); no average over
intercepts and no frozen-shift completion occurs.
\end{proposition}

\begin{proof}
The first arrow is
\eqref{eq:tpc58-terminal-grouped-equals-physical-high}.  Applying the
unsigned erasure and using
\eqref{eq:tpc58-literal-sign-calibration} gives, at a native coordinate
\(i\),
\[
 (\Eop A_H)_i
 =-\sum_{p>y}\sum_s\mu(s)b_{i,sp}
 =-\sum_{p>y}(v_p)_i,
\]
which is \eqref{eq:tpc58-fixed-shift-native-high-identity}.  All sums are
finite and retain the same literal carrier.
\end{proof}

\subsection{Raw norms and the second grouping gate}

The terminal atomic, residual grouped, and native-high energies on this
face are three different raw counting-measure quantities:
\begin{align}
 \cD_{h_0}^{>y}(G_X^L)
 &:=\sum_{\mathfrak u\in\cU_{h_0,>y}^{\rm term}}
      |X_{\mathfrak u}|^2,
 \label{eq:tpc58-very-high-terminal-diagonal}\\
 \cG_{h_0}^{>y}(G_X^L)
 &:=\sum_{i,s}
      \left|\sum_{\substack{\mathfrak u:\,
                    i(\mathfrak u)=i,\ s(\mathfrak u)=s}}
             X_{\mathfrak u}\right|^2
   =\|A_H\|_{\cH_{\rm res}^{H}}^2,
 \label{eq:tpc58-very-high-residual-grouped-energy}\\
 \cN_{H}^{>y}(h_0)
 &:=\sum_i\left|\sum_s A_H(i,s)\right|^2
   =\|\Eop A_H\|_{\cH_{\rm nat}}^2.
 \label{eq:tpc58-very-high-native-energy}
\end{align}
Equation \eqref{eq:tpc58-very-high-residual-grouped-energy} already forms
the coherent terminal-prime sum at fixed \((i,s)\).  The input to the native
Gram gate is therefore \(\cG_{h_0}^{>y}\), not
\(\cD_{h_0}^{>y}\).  The passage from
\eqref{eq:tpc58-very-high-residual-grouped-energy} to
\eqref{eq:tpc58-very-high-native-energy} is the second grouping
\((i,s)\to i\).

The singleton and phase-sector results of TPC--56--57 concern the first
grouping \(p\to(i,s)\) \citep{WangTPC56,WangTPC57}.  They do not constrain
the cross-\(s\) sum in
\eqref{eq:tpc58-very-high-native-energy}.  In particular:
\begin{enumerate}[label=\textup{(\roman*)},leftmargin=3em]
 \item a terminal-prime singleton at \((i,s)\) may coexist with many other
       singleton residual labels \(t\ne s\) in the same native coordinate;
 \item the adaptive phase sector of TPC--56 deletes terminal coordinates
       before the first grouping and is not the undeleted vector \(A_H\);
 \item the conditional undeleted floor of TPC--57 can provide a lower
       bound for \(\|A_H\|^2\), but no lower bound for
       \(\|\Eop A_H\|^2\) follows without a new residual-label comparison.
\end{enumerate}

Finally, \eqref{eq:tpc58-very-high-native-energy} is only the very-high
part of the native affine column on the declared direct slice.  Its complete
native affine energy contains
\begin{equation}
 v_0+\Eop A_H=v_0-\sum_{p\in\cP_y}v_p.
 \label{eq:tpc58-full-affine-sign-interface}
\end{equation}
Neither nonnegativity nor a triangle inequality permits a lower bound for
\(\|\Eop A_H\|^2\) to be returned to
\(\|v_0+\Eop A_H\|^2\).  Base completion is a separate affine gate.

The fixed-shift interface proved here is therefore exact but deliberately
limited: it aligns literal coefficients, signs, supports, and raw
coordinates through both grouping maps.  It proves no positive lower
singular value for \(\Eop\), no terminal activation, no residual-label
occupancy estimate, and no native affine lower bound.
